% =============================================================
%  Module Description — Networks (CS301-F28)
%  Shown on https://wosm.academy as the course page for CS301-F28.
%  Compile with: pdflatex cs301-module-description.tex
%  Write "TODO(owner): ..." where information is still missing;
%  the website highlights it.
% =============================================================
\documentclass[11pt,a4paper]{article}

\usepackage[utf8]{inputenc}
\usepackage[T1]{fontenc}
\usepackage{lmodern}
\usepackage[margin=2.2cm]{geometry}
\usepackage{booktabs}
\usepackage{tabularx}
\usepackage{array}
\usepackage{enumitem}
\usepackage{xcolor}
\usepackage{titlesec}
\usepackage{fancyhdr}
\usepackage[hidelinks]{hyperref}

% ---------- Module details (edit these) ----------------------
% \ModuleCode must equal the course's `code` in src/content/courses/.
\newcommand{\ModuleTitle}{Networks}
\newcommand{\ModuleCode}{CS301-F28}
\newcommand{\Programme}{Foundation programme}
\newcommand{\Level}{Semester 5}
\newcommand{\Semester}{Fall 2028}
\newcommand{\Credits}{6 ECTS}
\newcommand{\ModuleLeader}{Dr. Ali Khudiyev}
\newcommand{\Email}{ali.khudiyev@wosm.academy}
\newcommand{\OfficeHours}{Monday--Friday, 10:00--17:00 (Baku time, UTC+4)}
\newcommand{\Prerequisites}{CS204-S28 Operating Systems}

% ---------- Styling ------------------------------------------
\definecolor{accent}{RGB}{0,70,130}
\titleformat{\section}{\large\bfseries\color{accent}}{\thesection.}{0.5em}{}
\titlespacing*{\section}{0pt}{1.2em}{0.5em}
\setlist{itemsep=2pt, topsep=4pt}
\setlength{\parindent}{0pt}
\setlength{\parskip}{4pt}

\pagestyle{fancy}
\fancyhf{}
\lhead{\small \ModuleCode\ -- \ModuleTitle}
\rhead{\small Module Description}
\cfoot{\small \thepage}
\renewcommand{\headrulewidth}{0.4pt}

% =============================================================
\begin{document}

\begin{center}
  {\LARGE\bfseries\color{accent} \ModuleTitle}\\[4pt]
  {\large Module Description}
\end{center}
\vspace{6pt}

% ---------- Summary table (not shown on the website: the course page
% shows these details from the macros above) --------------------------
\begin{tabularx}{\textwidth}{>{\bfseries}l X >{\bfseries}l X}
\toprule
Module code   & \ModuleCode    & Credits      & \Credits \\
Programme     & \Programme     & Level        & \Level \\
Semester      & \Semester      & Prerequisites & \Prerequisites \\
Module leader & \ModuleLeader  & Email        & \href{mailto:\Email}{\Email} \\
Office hours  & \multicolumn{3}{l}{\OfficeHours} \\
\bottomrule
\end{tabularx}

% ---------- Sections (shown on the website) ------------------
\section{Module Overview}
This module explains how the Internet works, from the applications you use every day
down to the links that carry bits between machines. It follows the layered
architecture of the Internet from the top down: the application, transport, network
and link layers, with the problems each layer solves and the protocols that solve
them.

You will see why reliable communication over an unreliable network is possible, how
TCP shares bandwidth fairly, how routers find paths across the Internet, and how
local networks deliver frames. You will also write networked programs in C with the
POSIX sockets interface.
The module ends with the basics of network security.

\section{Aims}
\begin{itemize}
  \item Explain the architecture of the Internet and the principles behind its protocols.
  \item Show how key protocols at each layer work and why they were designed as they were.
  \item Give practical experience of writing and analysing networked programs.
  \item Introduce the security threats networks face and the basic defences against them.
\end{itemize}

\section{Learning Outcomes}
After completing this course, you will be able to:
\begin{enumerate}[label=LO\arabic*.]
  \item Explain layered network architectures and encapsulation, and calculate delay, loss and throughput in packet-switched networks.
  \item Explain application-layer protocols such as HTTP and DNS, and write client and server programs in C with the sockets interface.
  \item Explain and analyse transport-layer mechanisms: multiplexing, reliable data transfer, and TCP flow and congestion control.
  \item Explain IP addressing and forwarding, perform subnetting, and analyse link-state and distance-vector routing, OSPF and BGP.
  \item Explain link-layer functions: error detection, multiple access, Ethernet, switching and ARP.
  \item Explain the basic principles of network security, including cryptography, TLS and firewalls.
\end{enumerate}

\section{Syllabus}
\begin{tabularx}{\textwidth}{c X}
\toprule
\textbf{Unit} & \textbf{Topic} \\
\midrule
1  & Introduction: the Internet, packet switching; delay, loss and throughput \\
2  & Protocol layering and encapsulation \\
3  & Application layer I: HTTP and the web \\
4  & Application layer II: DNS, email, peer-to-peer and content distribution \\
5  & Socket programming in C: TCP and UDP clients and servers \\
6  & Transport layer: multiplexing, UDP, principles of reliable data transfer \\
7  & TCP: connection management and flow control \\
8  & Congestion control \\
9  & Network layer: forwarding, IPv4, addressing and subnetting, NAT, IPv6 \\
10 & Routing algorithms: link state and distance vector \\
11 & Routing in the Internet: OSPF, BGP; software-defined networking \\
12 & Link layer: error detection and correction, multiple access protocols \\
13 & Ethernet, switches, ARP and VLANs \\
14 & Wireless and mobile networks \\
15 & Network security: cryptography basics, TLS, firewalls \\
16 & Course Review \& Exam \\
\bottomrule
\end{tabularx}

\section{Teaching and Learning Methods}
The module is self-paced. All materials are provided on Moodle.
Each unit has a lecture and exercises. Practical tasks include writing socket
programs in C and capturing and analysing real network traffic with a packet analyser
such as Wireshark.

\begin{tabularx}{0.6\textwidth}{X r}
\toprule
\textbf{Activity} & \textbf{Hours} \\
\midrule
Lectures            & 32 \\
Labs / Tutorials    & 32 \\
Independent study   & 116 \\
\midrule
\textbf{Total}      & \textbf{180} \\
\bottomrule
\end{tabularx}

\section{Assessment}
Your knowledge will be assessed throughout the course with quizzes, practical
assignments and the final exam. The exam checks whether you can explain how
protocols work and solve quantitative problems on delay, throughput, addressing and
routing. The practical assignments check whether you can write and analyse networked
programs.

\begin{tabularx}{\textwidth}{X c c c}
\toprule
\textbf{Component} & \textbf{Weight} & \textbf{Due} & \textbf{LOs} \\
\midrule
Quizzes & 20\% & Throughout & LO1--LO6 \\
Practical assignments & 40\% & Throughout & LO2, LO3 \\
Final exam & 40\% & End of module & LO1--LO6 \\
\midrule
\textbf{Total} & \textbf{100\%} & & \\
\bottomrule
\end{tabularx}

\section{Reading List}
\textbf{Essential}
\begin{itemize}
  \item J.\ F.\ Kurose, K.\ W.\ Ross: \textit{Computer Networking: A Top-Down Approach}, 8th ed., Pearson, 2021.
\end{itemize}
\textbf{Recommended}
\begin{itemize}
  \item L.\ L.\ Peterson, B.\ S.\ Davie: \textit{Computer Networks: A Systems Approach}, 5th ed., Morgan Kaufmann, 2011. Updated edition free online at \url{https://book.systemsapproach.org}.
  \item A.\ S.\ Tanenbaum, N.\ Feamster, D.\ J.\ Wetherall: \textit{Computer Networks}, 6th ed., Pearson, 2021.
  \item K.\ R.\ Fall, W.\ R.\ Stevens: \textit{TCP/IP Illustrated, Volume 1: The Protocols}, 2nd ed., Addison-Wesley, 2011.
  \item W.\ R.\ Stevens, B.\ Fenner, A.\ M.\ Rudoff: \textit{UNIX Network Programming, Volume 1: The Sockets Networking API}, 3rd ed., Addison-Wesley, 2003.
\end{itemize}

\section{Materials and Licence}
Course materials are published under the
\href{https://creativecommons.org/licenses/by-sa/4.0/}{Creative Commons Attribution-ShareAlike 4.0}
licence unless stated otherwise. See \href{https://wosm.academy/licence/}{wosm.academy/licence} for details.

\section{Policies}
\begin{itemize}
  \item \textbf{Attendance:} Attendance is not mandatory. The module is self-paced, so you can study the materials when it suits you.
  \item \textbf{Late submissions:} Work submitted after the deadline is penalised.
  \item \textbf{Academic integrity:} All submitted work must be your own.
        See the \href{https://wosm.academy/code-of-conduct/}{Code of Conduct}.
\end{itemize}

\end{document}
